\documentclass[10pt,a4paper]{amsart}
\usepackage[margin=19mm]{geometry}
\usepackage{amsmath,amssymb,mathtools}
\usepackage[scr=rsfs,cal=euler]{mathalfa}
\usepackage{xfrac}
\usepackage[unicode,hidelinks,bookmarks=false]{hyperref}
\newcommand{\ie}{i.e.\ }
\newcommand{\cf}{cf.\ }
\newcommand{\resp}{respectively\ }
\newcommand{\Subsection}{\subsection}
\providecommand{\nobreakdash}{\nobreak\hbox{-}}
\setlength{\emergencystretch}{2em}
\multlinegap0pt
\usepackage{bm}
\usepackage{bbm}
\usepackage{dsfont}
\usepackage{xspace}
\usepackage{cancel}
\usepackage{mathtools}
\usepackage{amsthm}
\makeatletter
\@ifundefined{Theorem}{\newtheorem{Theorem}{Theorem}}{}
\@ifundefined{theorem}{\newtheorem{theorem}[Theorem]{Theorem}}{}
\makeatother
\title[A Greedy Version of the Frame Algorithm]{A Greedy Version of the Frame Algorithm}
\author{Brody Dylan Johnson}
\date{Source: arXiv:2506.18865v1 (CC BY 4.0)}
\begin{document}
\maketitle

\noindent\emph{Source and licence.} A Greedy Version of the Frame Algorithm. Version arXiv:2506.18865v1 (CC BY 4.0), distributed under the Creative Commons Attribution 4.0 International licence, as recorded in the arXiv licence metadata for this exact version and shown on the abstract page https://arxiv.org/abs/2506.18865v1. Full rights notice: licenses/arxiv-2506.18865-CC-BY-4.0.txt.\par\smallskip

\noindent\textit{Editorial scope.} Counted unit: the Theorem whose source label is \texttt{thm:robust} in the article (source0.tex lines 207--221, source label \texttt{thm:robust}), delivered together with the article's own complete original proof of it (source0.tex lines 223--266, one \texttt{proof} environment of 44 lines). No mathematics is written, repaired or re-derived: statement and proof are the article's own text line for line. Required context retained uncounted: thm:greedy (lines 132--146). Every definition, display and label of those lines is retained unchanged. Omitted from this package and not redistributed: 1--131, 147--206, 222--222, 267--360. Nothing inside a retained excerpt is omitted or altered: every retained body is delivered exactly as the article's lines are written, so no omission receipt of any kind accompanies this package. Citations: the retained excerpts carry no citation command at all, so no bibliography entry of the article is transcribed and no reference is redistributed. Reference adaptation: none was needed and none was made; every reference inside the delivered unit resolves inside it, so no unresolved reference or citation is produced. Printed slips kept literal: no printed word, symbol, hypothesis or number of the counted statement or of its proof was altered. Layout and class: the article's own class is \texttt{amsart} with packages including amssymb, latexsym, amsmath, amsthm, caption, graphicx, multicol, enumitem, fullpage; the wrapper is the lane's pinned \texttt{amsart} editorial wrapper at 10pt on A4, which restates the theorem-like environments the retained text needs and adds the article's own macro definitions that the retained text needs (none), with extra wrapper packages (bm, bbm, dsfont, xspace, cancel, mathtools). The delivered numbering is the unit's own. Assets: no retained span contains a figure environment, a graphics-inclusion command, a PDF-page inclusion, a table or a TikZ picture; no image file is copied into this package, the package declares no asset dependency and no image file is redistributed with it. Licence: version arXiv:2506.18865v1 of this article is distributed under the Creative Commons Attribution 4.0 International licence, as recorded in the arXiv licence metadata for this exact version and shown on the abstract page https://arxiv.org/abs/2506.18865v1. A full rights notice is delivered at licenses/arxiv-2506.18865-CC-BY-4.0.txt. Characters: every retained excerpt is pure ASCII, so the delivered engine, XeLaTeX with its default Latin Modern text font, reports no missing character.
\begin{theorem}[Greedy Frame Algorithm] \label{thm:greedy}
Let $\lbrace x_{j} \rbrace_{j\in J}$ be a frame for $\mathbb{H}$ with frame operator $S$.  Fix $x \in \mathbb{H}$.  Let $y_{0}=0$ and, for $n\ge 0$, define
\begin{equation} \label{eq:greedy}
\begin{aligned}
\alpha_{n} &= \frac{\Vert S(x-y_{n})\Vert^{2}}{\Vert S(x-y_{n})\Vert_{S}^{2}} \\[4pt]
y_{n+1} &= y_{n} + \alpha_{n} S(x-y_{n}).
\end{aligned}
\end{equation}

\noindent
Then, $y_{n}$ converges to $x$ as $n\rightarrow \infty$.  Moreover, if the optimal lower and upper frame bounds for $\lbrace x_{j} \rbrace_{j\in J}$ are $A$ and $B$, respectively, then
\begin{equation} \label{eq:greedy-bound}
\Vert x-y_{n} \Vert_{S} \le \big ( \tfrac{B-A}{B+A} \big )^{n} \Vert x \Vert_{S}.
\end{equation}
\end{theorem}

\begin{theorem} \label{thm:robust}
Let $\lbrace x_{j} \rbrace_{j\in J}$ be a frame for $\mathbb{H}$ with analysis operator $T$ and optimal lower and upper frame bounds $A$ and $B$, respectively.  Let $S$ represent the frame operator, $T^{*}T$.  Fix $x\in \mathbb{H}$ and $\delta_{0} >0$.  Let $c\in \ell^{2}(J)$ such that $\Vert Tx-c\Vert_{2} \le \delta_{0}$, define $y_{0}=0$ and, for $n\ge 0$, 
\begin{equation} \label{eq:greedy-2}
\begin{aligned}
\alpha_{n} &= \frac{\Vert T^{*}(c-Ty_{n})\Vert^{2}}{\Vert T^{*}(c-Ty_{n})\Vert_{S}^{2}} \\[4pt]
y_{n+1} &= y_{n} + \alpha_{n} T^{*}(c-Ty_{n}).
\end{aligned}
\end{equation}

\noindent
Then, 
\begin{equation} \label{eq:bound-2}
\Vert x-y_{n} \Vert_{S} \le \left ( \frac{B-A}{B+A} \right )^{n} \left ( \Vert x \Vert_{S} + 2\delta_{0}\right ) +  2\delta_{0}, \qquad n\ge 1.
\end{equation}
\end{theorem}

\begin{proof}
Fix $x\in \mathbb{H}$ and $c\in \ell^{2}(J)$ such that $\Vert Tx-c\Vert_{2} \le \delta_{0}$.  Define $\delta \in \ell^{2}(J)$ such that $c = Tx + \delta$ and observe that $\Vert \delta\Vert \le \delta_{0}$.  Let $P$ represent the orthogonal projection onto the range of $T$ and define $\tilde{x} \in \mathbb{H}$ so that $T\tilde{x}=Pc$, which implies that $c = T\tilde{x} + \tilde{c}$ where $\tilde{c}$ belongs to the null space of $T{^*}$.  With this setup, it follows from the second part of \eqref{eq:greedy-2} that
$$ \begin{aligned} 
    \tilde{x} - y_{n+1} &= \tilde{x} - y_{n} - \alpha_{n} T^{*}(c-T y_{n}) \\
                        &= \tilde{x} - y_{n} - \alpha_{n} T^{*}(T\tilde{x} + \tilde{c} - T y_{n}) \\
                        &= \tilde{x} - y_{n} - \alpha_{n} S(\tilde{x}-y_{n}) \\
                        &= (I-\alpha_{n}S) (\tilde{x}-y_{n}). 
    \end{aligned} $$

\noindent
In light of the previous calculation, one can write
$$ \Vert \tilde{x}-y_{n+1} \Vert_{S}^{2} = \Vert \tilde{x}-y_{n} \Vert_{S}^{2} - 2\alpha_{n} \langle S(\tilde{x}-y_{n}),\tilde{x}-y_{n}\rangle_{S} + \alpha_{n}^{2} \Vert S(\tilde{x}-y_{n}) \Vert_{S}^{2}. $$

\noindent
This expression for the error of approximation is minimized when
$$ \alpha_{n} = \frac{\Vert S(\tilde{x}-y_{n})\Vert^{2}}{\Vert S(\tilde{x}-y_{n})\Vert_{S}^{2}} = \frac{\Vert T^{*}(c-Ty_{n})\Vert^{2}}{\Vert T^{*}(c-Ty_{n})\Vert_{S}^{2}}$$

\noindent
and, repeating the argument of Theorem \ref{thm:greedy}, the error of approximation must satisfy
$$ \Vert \tilde{x}-y_{n+1} \Vert_{S} \le \left ( \frac{B-A}{B+A} \right ) \Vert \tilde{x}-y_{n} \Vert_{S}.$$

\noindent
It follows that
$$ \Vert \tilde{x}-y_{n} \Vert_{S} \le \left ( \frac{B-A}{B+A} \right )^{n} \Vert \tilde{x}-y_{0} \Vert_{S} = \left ( \frac{B-A}{B+A} \right )^{n} \Vert \tilde{x} \Vert_{S}, \qquad n \ge 1.$$

\noindent
It remains to relate this bound to the original vector $x$.  Since $T\tilde{x}=Pc$ is the orthogonal projection of $c$ onto the range of $T$, it follows that
$$ \Vert T\tilde{x} - c \Vert \le \Vert Tx - c \Vert \le \delta_{0}. $$

\noindent
In order to bound $\Vert x-\tilde{x}\Vert_{S}$, one can write
$$ \Vert x-\tilde{x} \Vert_{S}^{2} = \langle  S(x-\tilde{x}),  x-\tilde{x} \rangle = \Vert T(x-\tilde{x})\Vert^{2}, $$

\noindent
which shows that 
$$ \Vert x-\tilde{x}\Vert_{S} = \Vert T(x-\tilde{x})\Vert = \Vert (Tx-c)-(T\tilde{x}-c)\Vert \le 2\delta_{0}. $$ 

\noindent
Finally, combining this with the fact that $\Vert x-y_{n} \Vert_{S} \le \Vert \tilde{x}-y_{n}\Vert_{S} + \Vert x-\tilde{x}\Vert_{S}$ due to the triangle inequality, it follows that
$$ \Vert x-y_{n}\Vert_{S} \le \left ( \frac{B-A}{B+A} \right )^{n} \Vert \tilde{x} \Vert_{S} + 2 \delta_{0} \le \left ( \frac{B-A}{B+A} \right )^{n} \left ( \Vert x \Vert_{S} + 2\delta_{0} \right ) + 2 \delta_{0}, $$

\noindent
completing the proof of \eqref{eq:bound-2}.
\end{proof}

\end{document}
